\section*{Hoofdstuk \ref{ch:b2:biomolecules} --- Biomoleculen: aminozuren, suikers, lipiden en nucleotiden}

\begin{solution}{exo:b2:biomolecules:1}
\begin{itemize}
  \item Tussen pH 3 en 9 het \omterm{def:b2:biomolecules:amino-acid}{zwitterion} \ce{H3N+-CH(CH3)-COO-}.
  \item Bij pH 1 het kation \ce{H3N+-CH(CH3)-COOH}.
  \item Bij pH 12 het anion \ce{H2N-CH(CH3)-COO-}.
\end{itemize}
\end{solution}

\begin{solution}{exo:b2:biomolecules:2}
Glycine: $(2.35 + 9.78)/2 \approx 6.1$. Alanine: $(2.35 + 9.87)/2 \approx 6.1$.
% ledger: pka:glycine1, pka:glycine2, pka:alanine1, pka:alanine2
\end{solution}

\begin{solution}{exo:b2:biomolecules:3}
\ce{H2N-CH(CH3)-CO-NH-CH2-COOH}: alanine levert de N-terminus, glycine de C-terminus. Gly--Ala,
\ce{H2N-CH2-CO-NH-CH(CH3)-COOH}, is een constitutie-isomeer: de vrije aminogroep zit op glycine.
\end{solution}

\begin{solution}{exo:b2:biomolecules:4}
\ce{NH2} links: L. Prioriteiten \ce{NH2} > \ce{COOH} > \ce{CH2OH} > \ce{H}: L-serine is
(\textit{S}).
\end{solution}

\begin{solution}{exo:b2:biomolecules:5}
Groepen: N-terminaal ammonium (ongeveer 9), \omterm{def:b2:biomolecules:amino-acid}{zijketen} van lysine (10.7), \omterm{def:b2:biomolecules:amino-acid}{zijketen} van aspartaat (3.9),
C-terminale carbonzuurgroep (ongeveer 2). pH 1: $+1 +1 + 0 + 0 = +2$. pH 7: $+1 +1 -1 -1 = 0$. pH 12:
$0 + 0 - 1 - 1 = -2$ (de \omterm{def:b2:biomolecules:amino-acid}{zijketen} van lysine, 1.3 eenheden onder pH 12, is grotendeels neutraal).
\end{solution}

\begin{solution}{exo:b2:biomolecules:6}
Thionylchloride en warmte zouden ook andere groepen aanvallen, en vooral verliest het zeer reactieve \omterm{def:b2:acyl-substitution:derivative}{acylchloride}
van een aminozuur gemakkelijk zijn configuratie (het $\alpha$-waterstofatoom wordt zuur): het \omterm{def:b2:biomolecules:peptide}{peptide} zou
gedeeltelijk geracemiseerd zijn. Een \omterm{def:b2:biomolecules:peptide}{koppelingsreagens} activeert het zuur bij kamertemperatuur, onder milde omstandigheden.
\end{solution}

\begin{solution}{exo:b2:biomolecules:7}
O5 gebonden aan C2: een ring van C2, C3, C4, C5 en O, vijf atomen (een furanose). In \omterm{def:b2:biomolecules:anomer}{haworthprojectie}: ring-O
achteraan, C2 rechts met \ce{CH2OH} (C1) omlaag en \ce{OH} omhoog ($\beta$); op C3 \ce{OH} omhoog (links in
Fischer), op C4 \ce{OH} omlaag, C5 met \ce{CH2OH} (C6) omhoog.
\end{solution}

\begin{solution}{exo:b2:biomolecules:8}
$x_\alpha = (80.2 - 52.8)/(150.7 - 52.8) = 27.4/97.9 \approx 0.28$: 28~\% $\alpha$, 72~\% $\beta$.
\end{solution}

\begin{solution}{exo:b2:biomolecules:9}
Maltose houdt een vrij hemiacetaal op zijn tweede glucose, dat opengaat tot een aldehyde: reducerend. In sacharose
zijn beide \omterm{def:b2:biomolecules:anomer}{anomere koolstofatomen} in de \omterm{def:b2:biomolecules:glycoside}{glycosidische binding} betrokken: geen hemiacetaal, niet reducerend. Verdund zuur
hydrolyseert de \omterm{def:b2:biomolecules:glycoside}{glycosidische binding}, een acetaal: glucose en fructose komen vrij, en beide reduceren.
\end{solution}

\begin{solution}{exo:b2:biomolecules:10}
Trioleïne, \qty{885.45}{g/mol}, verbruikt drie \ce{KOH} (\qty{56.105}{g/mol}): $3 \times 56.105/885.45 =
0.190$~g per g, een verzepingsgetal van 190. Kortere ketens betekenen een kleinere molaire massa voor dezelfde
drie estergroepen: meer mol vet per gram, meer KOH.
% ledger: aw:C, aw:H, aw:O, aw:K
\end{solution}

\begin{solution}{exo:b2:biomolecules:11}
Bescherm het amine van glycine als \textit{tert}-butoxycarbonylcarbamaat, en het zuur van alanine als
methylester. Koppel met een carbodiimide. Verwijder het carbamaat met zuur en de ester met verdunde base
(orthogonale omstandigheden, zodat elk uiteinde afzonderlijk vrijgemaakt kan worden): Gly--Ala.
\end{solution}

\begin{solution}{exo:b2:biomolecules:12}
5$'$-ACGCAT-3$'$ (antiparallel gelezen: A paart met T, G met C). Paren: A--T, T--A, G--C, C--G, G--C, T--A:
$3 \times 2 + 3 \times 3 = 15$ waterstofbruggen.
\end{solution}

\begin{solution}{pb:b2:biomolecules:1}
\textbf{1.} L-Asparaginezuur, L-fenylalanine en methanol.
\textbf{2.} In beide \ce{COOH} bovenaan, de \omterm{def:b2:biomolecules:amino-acid}{zijketen} onderaan (\ce{CH2COOH}, \ce{CH2C6H5}), \ce{NH2}
links.
\textbf{3.} (\textit{S}) voor beide.
\textbf{4.} Twee stereocentra: vier stereo-isomeren.
\textbf{5.} \ce{H2N-CH(CH2COOH)-CO-NH-CH(CH2C6H5)-COOCH3}: de \omterm{def:b2:biomolecules:peptide}{peptidebinding} verbindt de twee
residuen; de ester is de \ce{COOCH3} aan de C-terminus.
\textbf{6.} Smaakreceptoren zijn chiraal: ze binden één stereo-isomeer en niet zijn spiegelbeeld of zijn
diastereo-isomeren, zoals de inleiding zei.
\textbf{7.} De N-terminale \ce{NH3+} en de \ce{COOH} van de \omterm{def:b2:biomolecules:amino-acid}{zijketen} van asparaginezuur; het C-terminale zuur
is een methylester, zonder zuur waterstofatoom.
\textbf{8.} Bij pH 2: \ce{NH3+} ($+1$), \ce{COOH} van de \omterm{def:b2:biomolecules:amino-acid}{zijketen} (0): $+1$.
\textbf{9.} Bij pH 7: $+1 - 1 = 0$.
\textbf{10.} Bij pH 12: $0 - 1 = -1$.
\textbf{11.} Tussen de twee $\mathrm pK_a$ van de groepen die veranderen: $(3.9 + 9.9)/2 \approx 6.9$.
\textbf{12.} In het dipeptide draagt het koolstofatoom van het amine een amide in plaats van een carboxylaat: de
naburige groepen en ladingen veranderen, en daarmee ook de $\mathrm pK_a$ (het N-terminale ammonium van een
\omterm{def:b2:biomolecules:peptide}{peptide} is duidelijk zuurder dan dat van een vrij aminozuur). Het model geeft de vorm, niet de
exacte waarden.
\textbf{13.} Zuur en amine vormen eerst een zout; verhitten zou mengsels geven (Asp--Asp, langere ketens,
de verkeerde carbonzuurgroep) en de stereocentra racemiseren.
\textbf{14.} Zijn aminogroep en de carbonzuurgroep van zijn \omterm{def:b2:biomolecules:amino-acid}{zijketen}.
\textbf{15.} Het maakt van de \ce{OH} van het zuur een goede vertrekkende groep, een geactiveerd derivaat dat het
amine bij kamertemperatuur aanvalt (acylsubstitutie).
\textbf{16.} Het $\beta$-dipeptide, verbonden via de \omterm{def:b2:biomolecules:amino-acid}{zijketen}: een isomeer van aspartaam, geen aspartaam.
\textbf{17.} Het amine als benzyloxycarbonylcarbamaat en de \omterm{def:b2:biomolecules:amino-acid}{zijketen} als benzylester: beide worden
samen verwijderd door hydrogenolyse over palladium, die de methylester ongemoeid laat (orthogonaal
daarop).
\textbf{18.} Een geactiveerd zuur heeft een zuur $\alpha$-waterstofatoom; base haalt het weg en het stereocentrum
racemiseert.
\textbf{19.} Bescherm asparaginezuur (carbamaat op N, benzylester op de \omterm{def:b2:biomolecules:amino-acid}{zijketen}); activeer de vrije
$\alpha$-\ce{COOH} met het \omterm{def:b2:biomolecules:peptide}{koppelingsreagens}; voeg de methylester van L-fenylalanine toe; hydrogenolyseer.
\textbf{20.} De methylester en de \omterm{def:b2:biomolecules:peptide}{peptidebinding}; de ester eerst, want een amide is veel minder reactief.
\textbf{21.} Het dipeptide Asp--Phe (vrij zuur) en methanol.
\textbf{22.} Het zoete molecuul wordt verbruikt: zijn hydrolyseproducten nemen zijn plaats niet in.
\textbf{23.} $14 \times 12.011 + 18 \times 1.008 + 2 \times 14.007 + 5 \times 15.999 =
\qty{294.307}{g/mol}$.
\textbf{24.} $0.180/294.307 = \qty{6.116e-4}{mol}$, \qty{0.612}{mmol}.
\textbf{25.} Eén methanol per aspartaam: $0.6116 \times 32.042 = \boldsymbol{\approx \qty{19.6}{mg}}$
methanol.
\end{solution}
